\documentclass[aspectratio=169,10pt]{beamer}
\usepackage{cstheme}
% ===========================================================================
% Repeated figures as macros: the same picture is literally the same drawing
% wherever it appears, and cannot drift. Page millimetres, origin top-left.
% ===========================================================================
\newcommand{\csx}{\tikz[baseline]{\draw[csRed,line width=1.05pt,line cap=round]
  (0,0.2mm)--(2.4mm,2.4mm) (0,2.4mm)--(2.4mm,0.2mm);}}

% A subgraph body drawn as a dashed box.  \bodybox{x}{y}{w}{h}{label}
\newcommand{\bodybox}[5]{%
  \node[csbody,anchor=north west,minimum width=#3 mm,minimum height=#4 mm,fill=white]
       at (#1,#2) {};
  \node[cslabel,anchor=north west,inner sep=0pt,align=left] at (#1+1.6,#2+1.2) {#5};}

% A row of port labels down the left or right edge of a body box.
% \ports{x}{y0}{dy}{style}{a/b/c}
\newcommand{\ports}[5]{%
  \foreach \p [count=\i from 0] in {#5}
    {\node[#4,anchor=east,inner sep=0pt] at (#1,{#2+\i*#3}) {\p};}}
\newcommand{\portsr}[5]{%
  \foreach \p [count=\i from 0] in {#5}
    {\node[#4,anchor=west,inner sep=0pt] at (#1,{#2+\i*#3}) {\p};}}

% A Scan/Loop/If node drawn large enough to hold a body.
% \cfnode{x}{y}{w}{h}{title}{style}
\newcommand{\cfnode}[6]{%
  \node[#6,anchor=north west,minimum width=#3 mm,minimum height=#4 mm,align=left]
       at (#1,#2) {};
  \node[anchor=north west,inner sep=0pt,font=\latosemi\footnotesize,text=csDeep]
       at (#1+2.4,#2+2.2) {#5};}


\title{Differentiating control flow}
\subtitle{A strategy for \texttt{If}, \texttt{Scan} and \texttt{Loop}}
\author{onnx-ad}
\date{22 September 2026}
\setcstitlechips{\cstitlechip{onnx-ad 0.1.0}\;\cstitlechip{design brief}}
\setcsfootlabel{onnx-ad \textemdash{} differentiating control flow}

\begin{document}
\csmaketitle

% ========================================================== 1. vocabulary ==
\begin{frame}{One graph, three readings}
\cscanvas{%
  \node[cslabel,anchor=west,inner sep=0pt] at (9,34) {\textsc{primal}};
  \node[cslabel,anchor=west,inner sep=0pt] at (9,52) {\textsc{forward adds}};
  \node[cslabel,anchor=west,inner sep=0pt] at (9,64) {\textsc{reverse adds}};

  % --- primal band, the grid every other band aligns to
  \node[csghost,anchor=north west,minimum width=22mm] (pm) at (41,30) {MatMul};
  \node[csghost,anchor=north west,minimum width=22mm] (pt) at (77,30) {Tanh};
  \draw[csthin] (33,34) -- (pm.west);  \node[cslabel,anchor=south east] at (33,34.8) {$x$};
  \draw[csthin] (pm.east) -- (pt.west); \node[cslabel,anchor=south] at (70,34.8) {$h$};
  \draw[csthin] (pt.east) -- (111,34);  \node[cslabel,anchor=south west] at (111,34.8) {$y$};

  % --- forward band
  \node[cshot,anchor=north west,minimum width=22mm] (fm) at (41,48) {MatMul};
  \node[cshot,anchor=north west,minimum width=22mm] (ft) at (77,48) {$\times(1-y^{2})$};
  \draw[csflow] (33,52) -- (fm.west);
  \node[cslabel,anchor=south east] at (33,52.8) {\texttt{fwd\_x}};
  \draw[csflow] (fm.east) -- (ft.west);
  \draw[csflow] (ft.east) -- (111,52);
  \node[cslabel,anchor=south west] at (111,52.8) {\texttt{fwd\_y}};
  \draw[cslead] (pt.south) -- (ft.north);

  % --- reverse band, read right to left
  \node[cshot,anchor=north west,minimum width=22mm,draw=csRed,fill=csRed!5,text=csRed!60!black]
       (rm) at (41,60) {$\cdot\,W^{\!\top}$};
  \node[cshot,anchor=north west,minimum width=22mm,draw=csRed,fill=csRed!5,text=csRed!60!black]
       (rt) at (77,60) {$\times(1-y^{2})$};
  \draw[csback] (111,64) -- (rt.east);
  \node[cslabel,anchor=south west] at (111,64.8) {\texttt{adj\_y}};
  \draw[csback] (rt.west) -- (rm.east);
  \draw[csback] (rm.west) -- (33,64);
  \node[cslabel,anchor=south east] at (33,64.8) {\texttt{adj\_x}};
  \draw[cslead] (pt.south) .. controls (99,50) and (99,54) .. (rt.north);
}
\begin{csbox}{140mm}{9mm}{70mm}
\muted{\footnotesize A pass \emph{adds} a band and keeps the primal, so a nonlinear rule
reads the tensor the primal already produced \textemdash{} the dashed links. Everything
after this slide is the same picture with a subgraph inside a node.}
\end{csbox}
\end{frame}

% ============================================================== 2. order ===
\begin{frame}{Why \texttt{If}, then \texttt{Scan}, then \texttt{Loop}}
\cscanvas{%
  \node[cslabel,anchor=west,inner sep=0pt] at (44,27) {\textsc{new capability needed}};
  \node[cslabel,anchor=west,inner sep=0pt] at (100,27) {\textsc{reverse mode needs}};
  \foreach \n/\y/\a/\b in {%
    \texttt{If}/30/{recursion into subgraphs,\\outer-scope capture}/{nothing iterative},
    \texttt{Scan}/43/{a tape, and a sweep in\\the opposite direction}/{a trip count known
      \textbf{before}\\the primal runs},
    \texttt{Loop}/56/{a trip count known only\\\textbf{after} the primal runs}/{the \texttt{Scan}
      machinery,\\fed by that trip count}}{%
    \node[cshot,anchor=north west,minimum width=24mm,minimum height=9mm] at (9,\y) {\n};
    \node[anchor=north west,text width=52mm,align=left,inner sep=0pt,
          font=\footnotesize,text=csInk] at (44,\y+1.2) {\a};
    \node[anchor=north west,text width=50mm,align=left,inner sep=0pt,
          font=\footnotesize,text=csInk] at (100,\y+1.2) {\b};
    \draw[cslead] (9,\y+11.5) -- (151,\y+11.5);}
}
\begin{csbox}{140mm}{9mm}{71mm}
\begin{cstake}
The adjoint of a \texttt{Loop} is a \texttt{Scan}: once the primal loop has run, its trip
count is data.
\end{cstake}
\end{csbox}
\end{frame}

% ================================================================= 3. If ===
\begin{frame}{\texttt{If} \textemdash{} differentiate inside the branch}
\cscanvas{%
  \foreach \x/\lbl in {9/{\textsc{primal}}, 57/{\textsc{forward}}, 105/{\textsc{reverse}}}
    {\node[cslabel,anchor=west,inner sep=0pt] at (\x,27) {\lbl};}

  % ---- primal
  \cfnode{9}{30}{38}{32}{If}{csghost}
  \bodybox{13}{39}{30}{9}{then}
  \bodybox{13}{50}{30}{9}{else}
  \node[csport,anchor=west,inner sep=0pt] at (13,35.8) {\texttt{cond}};
  \draw[csthin] (47,44) -- (53,44); \node[csport,anchor=east,inner sep=0pt] at (53,41.8) {$y$};

  % ---- forward
  \cfnode{57}{30}{38}{32}{If}{cshot}
  \bodybox{61}{39}{30}{9}{then $+$ tangent}
  \bodybox{61}{50}{30}{9}{else $+$ tangent}
  \node[csport,anchor=west,inner sep=0pt] at (61,35.8) {\texttt{cond}};
  \draw[csthin] (95,44) -- (101,44); \node[csport,anchor=east,inner sep=0pt] at (101,41.8) {$y$};
  \draw[csflow] (95,54) -- (101,54);
  \node[csportt,anchor=east,inner sep=0pt] at (101,51.8) {\texttt{fwd\_y}};

  % ---- reverse
  \cfnode{105}{30}{38}{32}{If}{cshot}
  \bodybox{109}{39}{30}{9}{then adjoint}
  \bodybox{109}{50}{30}{9}{else adjoint}
  \draw[csback] (151,44) -- (143,44);
  \node[csport,anchor=east,inner sep=0pt] at (151,41.8) {\texttt{adj\_y}};
  \draw[csback] (143,54) -- (151,54);
  \node[csport,anchor=east,inner sep=0pt] at (151,51.8) {\texttt{adj\_}$w$};

  \foreach \x/\t in {9/{branches have no inputs:\\everything is captured},
                     57/{each branch output grows\\by its own tangent},
                     105/{one output per \emph{captured}\\outer value}}
    {\node[cslabel,anchor=north west,text width=42mm,inner sep=0pt] at (\x,64) {\t};}
}
\begin{csbox}{140mm}{9mm}{70mm}
\begin{cstake}
Only the taken branch computes its derivative \textemdash{} which is why not to evaluate both
and select with \texttt{Where}.
\end{cstake}
\end{csbox}
\end{frame}

% ====================================================== 4. Scan, forward ===
\begin{frame}{\texttt{Scan} forward \textemdash{} one \texttt{Scan}, no tape}
\cscanvas{%
  \node[cslabel,anchor=west,inner sep=0pt] at (9,27) {\textsc{primal}};
  \cfnode{9}{31}{60}{22}{Scan}{csghost}
  \bodybox{25}{37}{28}{12}{body}
  \ports{23.5}{41}{5}{csport}{$s$,$x$}
  \portsr{54.5}{41}{5}{csport}{$s'$,$y$}
  \node[cslabel,anchor=north west,inner sep=0pt] at (9,55) {\texttt{num\_scan\_inputs} $=M$};

  \node[cslabel,anchor=west,inner sep=0pt] at (83,27) {\textsc{forward}};
  \cfnode{83}{31}{60}{34}{Scan}{cshot}
  \bodybox{99}{37}{28}{24}{body}
  \ports{97.5}{41}{5}{csport}{$s$,$x$}
  \ports{97.5}{51}{5}{csportt}{\texttt{t\_}$s$,\texttt{t\_}$x$}
  \portsr{128.5}{41}{5}{csport}{$s'$,$y$}
  \portsr{128.5}{51}{5}{csportt}{\texttt{t\_}$s'$,\texttt{t\_}$y$}
  \node[cslabel,anchor=north west,inner sep=0pt] at (83,67) {\texttt{num\_scan\_inputs} $=2M$};
  \draw[csflow] (71,45) -- (80,45);
}
\begin{csbox}{140mm}{9mm}{71mm}
\begin{cstake}
The tangent of a scan is a scan of the same length whose state is the pair $(s,\ \dot s)$:
two\texttimes{} the body, any number of seeds, nothing stored.
\end{cstake}
\end{csbox}
\end{frame}

% ====================================================== 5. Scan, reverse ===
\begin{frame}{\texttt{Scan} reverse \textemdash{} tape the state, sweep backwards}
\cscanvas{%
  \cfnode{9}{29}{62}{26}{Scan \textemdash{} primal $+$ tape}{csghost}
  \bodybox{26}{37}{28}{14}{body}
  \ports{24.5}{42}{5}{csport}{$s$,$x$}
  \portsr{55.5}{42}{5}{csport}{$s'$,$y$}
  \node[cstape,anchor=north,minimum width=32mm] (tp) at (40,59) {tape: $s$ per iteration};
  \draw[cstapeflow] (40,55) -- (40,59);

  \cfnode{85}{29}{62}{30}{Scan \textemdash{} reversed}{cshot}
  \bodybox{104}{37}{30}{18}{recompute,\\then adjoint}
  \ports{102.5}{42}{5}{csport}{$\leftarrow$\,tape,$\leftarrow x$,$\leftarrow$\,\texttt{adj\_}$y$}
  \portsr{135.5}{42}{5}{csport}{\texttt{adj\_}$s$,\texttt{adj\_}$x$,\texttt{acc\_}$w$}
  \draw[cstapeflow] (tp.east) .. controls (72,63) and (78,50) .. (85,46);
  \node[cslabel,anchor=north west,inner sep=0pt] at (85,61)
    {\texttt{direction\,=\,1} on every scan input and output};
}
\begin{csbox}{140mm}{9mm}{69mm}
\begin{cstake}
The taped input state is the minimal recomputation root \textemdash{} and a captured weight
accumulates in the \emph{state}, so its adjoint costs $O(|w|)$, not $O(T\!\cdot\!|w|)$.
\end{cstake}
\end{csbox}
\end{frame}

% ============================================================== 6. Loop ====
\begin{frame}{\texttt{Loop} reverse \textemdash{} the trip count becomes data}
\cscanvas{%
  \cfnode{9}{29}{44}{22}{Loop \textemdash{} primal $+$ tape}{csghost}
  \bodybox{20}{37}{26}{11}{body}
  \node[cstape,anchor=north,minimum width=24mm] (tp) at (31,56) {tape};
  \draw[cstapeflow] (31,51) -- (31,56);
  \node[cshot,anchor=north west,minimum width=34mm] (sh) at (62,56)
       {\texttt{Shape(tape)[0]}\,$=T$};
  \draw[cstapeflow] (tp.east) -- (sh.west);

  \cfnode{62}{27}{84}{24}{If\ \ $T>0$}{cshot}
  \bodybox{70}{34}{70}{10}{then: the reversed \texttt{Scan} from the previous slide}
  \node[cslabel,anchor=north west,text width=70mm,inner sep=0pt] at (70,46)
    {else: the zero-initialised accumulators};
  \draw[csflow] (sh.north) -- (79,51);
}
\begin{csbox}{140mm}{9mm}{63mm}
\muted{\footnotesize A reversed \texttt{Loop} tolerates $T=0$ natively, but needs
\texttt{Gather(tape,\,T-1-i)} \emph{and} a final flip of its scan outputs.}
\end{csbox}
\begin{csbox}{140mm}{9mm}{71mm}
\begin{cstake}
The reverse sweep is a \texttt{Scan}, not a \texttt{Loop}: reading backwards removes all index
arithmetic.
\end{cstake}
\end{csbox}
\end{frame}

% ======================================================= 7. checked + plan =
\begin{frame}{Checked against the runtime, and the order of work}
\cscanvas{%
  \node[cslabel,anchor=west,inner sep=0pt] at (9,26) {\textsc{ONNX Runtime 1.30 \textemdash{} measured, not assumed}};
  \foreach \m/\y/\t in {%
    \cscheck/33/{outer-scope capture, two levels deep},
    \cscheck/40/{reversed scan input \emph{and} output directions},
    \cscheck/47/{a scan state carrying a trailing seed axis},
    \cscheck/54/{loop trip count from \texttt{Shape(tape)[0]}},
    \csx/61/{\textbf{a zero-length \texttt{Scan} is rejected}},
    \cscheck/68/{an \texttt{If} guarding that empty case}}{%
    \node[anchor=west,inner sep=0pt] at (9,\y) {\m};
    \node[anchor=west,inner sep=0pt,font=\footnotesize,text=csInk,text width=56mm]
         at (14,\y) {\t};}

  \node[cslabel,anchor=west,inner sep=0pt] at (84,26) {\textsc{staging}};
  \foreach \n/\y/\t in {%
    0/34/{\texttt{unroll} for a static trip count \textemdash{} small, useful
          alone, and the test oracle},
    1/46/{scope chain, shared name allocator, capture analysis; then \texttt{If}},
    2/58/{\texttt{Scan} forward, then \texttt{Scan} reverse},
    3/68/{\texttt{Loop}, on top of \texttt{Scan}'s}}{%
    \node[cshot,anchor=north west,minimum width=7mm,minimum height=7mm,
          font=\latosemi\footnotesize] at (84,\y-3) {\n};
    \node[anchor=north west,inner sep=0pt,font=\footnotesize,text=csInk,text width=52mm]
         at (95,\y-3) {\t};}
}
\begin{csbox}{140mm}{9mm}{73mm}
\begin{cstake}
Unrolling is the oracle: two Jacobians from two independent implementations, needing no
framework to referee them.
\end{cstake}
\end{csbox}
\end{frame}

\end{document}
